\begin{ouverture}[geometrie]{15}{Angles et bissectrices}{Partie 3 · Espace et géométrie}
{Vocabulaire · Rapporteur · Angles associés · Bissectrice}
\objectif{Nommer, mesurer et construire un angle au rapporteur.}
\objectif{Utiliser angles adjacents, opposés par le sommet et supplémentaires.}
\objectif{Définir et construire la bissectrice d’un angle.}
\prerequis{Demi-droites ; angle droit ; lecture d’une graduation.}
\end{ouverture}
\pagestyle{otMain}\essentiel\fontsize{11}{15}\selectfont
\phantomsection\label{cours-15-1}
\essSection{1}{Nommer et reconnaître les angles}
\begin{definition}Un angle est formé de deux demi-droites de même origine : ses côtés. Cette origine est son sommet. L’angle formé par les demi-droites [OA) et [OB) se note $\widehat{AOB}$ : la lettre du sommet est au milieu. Les angles se mesurent en degrés.\end{definition}
\begin{definition}Un angle est nul s’il mesure $0^\circ$ ; aigu s’il mesure entre $0^\circ$ et $90^\circ$ ; droit s’il mesure $90^\circ$ ; obtus entre $90^\circ$ et $180^\circ$ ; plat s’il mesure $180^\circ$ ; plein s’il mesure $360^\circ$.\end{definition}
\essSection{2}{Mesurer et construire au rapporteur}
\begin{savoirfaire}[Mesurer un angle]On place le centre du rapporteur sur le sommet et une graduation 0 sur un côté. On lit la mesure à l’endroit où passe le second côté, sur la graduation qui part de 0.\end{savoirfaire}
\label{fig-cours-15-1-cours-rapporteur}\figureGeo{cours-rapporteur}

\begin{savoirfaire}[Construire un angle]Pour construire un angle $\widehat{xAy}$ de $65^\circ$ : tracer [Ax), placer le centre du rapporteur en A et le 0 sur [Ax), marquer un point en face de 65, puis tracer la demi-droite [Ay).\end{savoirfaire}
\clearpage
\phantomsection\label{cours-15-2}
\essSection{3}{Angles adjacents, opposés et supplémentaires}
\begin{definition}Deux angles adjacents ont le même sommet et un côté commun, et sont situés de part et d’autre de ce côté. Deux angles sont supplémentaires si la somme de leurs mesures vaut $180^\circ$.\end{definition}
\begin{propriete}Quand deux droites se coupent, deux angles opposés par le sommet ont la même mesure. Deux angles adjacents qui forment un angle plat sont supplémentaires.\end{propriete}
\label{fig-cours-15-2-cours-paires-bissectrice}\figureGeo{cours-paires-bissectrice}

\essSection{4}{Définir et construire une bissectrice}
\begin{definition}La bissectrice d’un angle saillant, c’est-à-dire de mesure inférieure à $180^\circ$, est la demi-droite issue du sommet qui partage cet angle en deux angles adjacents de même mesure.\end{definition}
\begin{savoirfaire}[Au compas]Un arc de centre O coupe les côtés en M et N. Deux arcs de même rayon, de centres M et N, se coupent en P. La demi-droite [OP) est la bissectrice. Au rapporteur, on marque simplement la moitié de la mesure.\end{savoirfaire}

\clearpage\rubrique{Découverte}
\begin{blocLivre}{ex-15-1}\begin{exercice}[1]{Nommer un angle}
Nommer l’angle marqué en vert, puis donner son sommet et ses côtés.
\label{fig-ex-15-1}\figureGeo{ang-1}
\end{exercice}\end{blocLivre}
\begin{blocLivre}{ex-15-2}\begin{exercice}[1]{Aigu, droit, obtus…}
Pour chacun des angles a, b, c, d et e, dire s’il est aigu, droit, obtus ou plat.
\label{fig-ex-15-2}\figureGeo{ang-2}
\end{exercice}\end{blocLivre}
\begin{blocLivre}{ex-15-3}\begin{exercice}[1]{Lire le rapporteur}
Lire la mesure de l’angle $\widehat{AOB}$.
\label{fig-ex-15-3}\figureGeo{ang-3}
\end{exercice}\end{blocLivre}
\begin{blocLivre}{ex-15-4}\begin{exercice}[1]{Avec l’autre graduation}
Lire la mesure de l’angle $\widehat{COD}$.
\label{fig-ex-15-4}\figureGeo{ang-4}
\end{exercice}\end{blocLivre}
\begin{blocLivre}{ex-15-5}\begin{exercice}[1]{Estimer}
Sans rapporteur, associer chacun des angles a, b et c à l’une des mesures $45^\circ$, $100^\circ$ et $160^\circ$.
\label{fig-ex-15-5}\figureGeo{ang-5}
\end{exercice}\end{blocLivre}

\rubrique{Entraînement}
\begin{blocLivre}{ex-15-6}\begin{exercice}[2]{Construire un angle}
Reproduire la demi-droite [Ax), puis construire la demi-droite [Ay) telle que $\widehat{xAy}=65^\circ$.
\label{fig-ex-15-6}\figureGeo{ang-6}
\end{exercice}\end{blocLivre}
\begin{blocLivre}{ex-15-7}\begin{exercice}[2]{Adjacents ou non ?}
Dans chacune des trois figures, les deux angles marqués sont-ils adjacents ? Justifier.
\label{fig-ex-15-7}\figureGeo{ang-7}
\end{exercice}\end{blocLivre}
\begin{blocLivre}{ex-15-8}\begin{exercice}[2]{Supplémentaires}
Les points A, O et B sont alignés. L’angle $\widehat{AOC}$ mesure $128^\circ$. Calculer la mesure de l’angle $\widehat{COB}$.
\label{fig-ex-15-8}\figureGeo{ang-8}
\end{exercice}\end{blocLivre}
\begin{blocLivre}{ex-15-9}\begin{exercice}[2]{Opposés par le sommet}
Deux droites se coupent en O. L’un des angles mesure $47^\circ$. Calculer les trois autres.
\label{fig-ex-15-9}\figureGeo{ang-9}
\end{exercice}\end{blocLivre}
\begin{blocLivre}{ex-15-10}\begin{exercice}[2]{La bissectrice au rapporteur}
L’angle $\widehat{xOy}$ mesure $84^\circ$. Construire sa bissectrice au rapporteur et donner la mesure des deux angles obtenus.
\label{fig-ex-15-10}\figureGeo{ang-10}
\end{exercice}\end{blocLivre}

\rubrique{Problèmes}
\begin{blocLivre}{ex-15-11}\begin{exercice}[3]{Autour d’un point}
Quatre angles adjacents font le tour du point O. Trois d’entre eux mesurent $110^\circ$, $95^\circ$ et $80^\circ$. Calculer la mesure du quatrième.
\label{fig-ex-15-11}\figureGeo{ang-11}
\end{exercice}\end{blocLivre}
\begin{blocLivre}{ex-15-12}\begin{exercice}[3]{Les aiguilles}
Le cadran d’une horloge est partagé en 12 heures. Quel angle forment les deux aiguilles à 3 h ? À 4 h, comme sur le dessin ? À 5 h ?
\label{fig-ex-15-12}\figureGeo{ang-12}
\end{exercice}\end{blocLivre}
\begin{blocLivre}{ex-15-13}\begin{exercice}[3]{La bissectrice au compas}
Reproduire l’angle $\widehat{xOy}$ de $70^\circ$, puis construire sa bissectrice au compas, sans rapporteur. Quelle est la mesure de chaque moitié ?
\label{fig-ex-15-13}\figureGeo{ang-13}
\end{exercice}\end{blocLivre}
\begin{blocLivre}{ex-15-14}\begin{exercice}[3]{Deux bissectrices}
Les points A, O et B sont alignés et $\widehat{AOC}=130^\circ$. On trace les bissectrices des angles $\widehat{AOC}$ et $\widehat{COB}$. Quel angle forment-elles ?
\label{fig-ex-15-14}\figureGeo{ang-14}
\end{exercice}\end{blocLivre}
\begin{blocLivre}{ex-15-15}\begin{exercice}[3]{Vrai ou faux}
Inès affirme : « L’angle a est plus grand que l’angle b, car ses côtés sont plus longs. » A-t-elle raison ?
\label{fig-ex-15-15}\figureGeo{ang-15}
\end{exercice}\end{blocLivre}

\clearpage
